\chapter{Model Customization -- Fine-Tuning und Adaptation}
\label{chap:finetuning}

% ============================================================
% LERNZIELE
% ============================================================
\begin{lernziele}
Nach diesem Kapitel kannst du:
\begin{itemize}
    \item Die Entscheidungsmatrix RAG vs Fine-Tuning vs Prompt Engineering anwenden
    \item LoRA und QLoRA korrekt konfigurieren (Rank, Alpha, Target-Module)
    \item Ein Instruction-Dataset für Fine-Tuning aufbereiten und validieren
    \item Einen Fine-Tuning-Job auf GPU-Infrastruktur planen und kostenschätzen
    \item Ein fine-getuntes Modell evaluieren (catastrophic forgetting erkennen)
    \item Multi-LoRA-Serving mit vLLM konfigurieren
\end{itemize}
\end{lernziele}

% ============================================================
% MOTIVATION
% ============================================================
\section{Motivation}

RAG und Prompt Engineering decken viele Anwendungsfälle ab. Aber sie haben Grenzen: Ein Modell, das in einem spezifischen Format antworten, eine bestimmte Terminologie verwenden oder einen speziellen Stil generieren muss, lässt sich durch Prompts allein nur schwer zuverlässig steuern. Und manche Domänen -- Medizin, Jura, Finanzen -- erfordern Wissen, das selbst das beste Prompt-Design nicht aus dem Basismodell extrahieren kann.

Fine-Tuning schliesst diese Lücke: Es passt die Modellgewichte an deine spezifische Domäne an. Ein auf medizinischen Texten fine-getuntes Modell versteht klinische Terminologie. Ein auf Unternehmensdokumente getuntes Modell schreibt im Corporate Voice. Ein auf SQL fine-getuntes Modell generiert korrekte Datenbankabfragen.

Der Haken: Fine-Tuning ist teuer, aufwändig und risikoreich. Falsch durchgeführt, verliert das Modell seine allgemeinen Fähigkeiten (catastrophic forgetting). Dieses Kapitel zeigt, wann Fine-Tuning sinnvoll ist und wie du es korrekt durchführst.

% ============================================================
% GRUNDLAGEN
% ============================================================
\section{Grundlagen -- RAG vs Fine-Tuning vs Prompt Engineering}

Die drei Techniken schliessen sich nicht aus, sondern ergänzen sich:

\begin{table}[H]
\centering
\begin{tabularx}{\linewidth}{lXXXX}
\toprule
\textbf{Kriterium} & \textbf{Prompt Engineering} & \textbf{RAG} & \textbf{Fine-Tuning} \\
\midrule
Aufwand & Gering (Stunden) & Mittel (Tage) & Hoch (Wochen) \\
Kosten & Null bis gering & API-Kosten + Vector-Infrastruktur & GPU-Kosten + Dataset-Erstellung \\
Flexibilität & Sofort änderbar & Daten-Update genügt & Neues Training nötig \\
Domänenwissen & Keines (Prompt lenkt) & Aus Dokumenten extrahiert & In Gewichte integriert \\
Output-Format & Prompt-abhängig, nicht erzwingbar & Prompt-abhängig & Format erlernbar \\
Wissensaktualität & Aktuell durch Prompt & Aktuell durch Index & Zum Trainingszeitpunkt \\
Qualität bei Spezialisierung & Niedrig bis mittel & Mittel bis hoch & Hoch \\
\bottomrule
\end{tabularx}
\caption{RAG vs Fine-Tuning vs Prompt Engineering}
\label{tab:ft-vs-rag}
\end{table}

\paragraph{Entscheidungsmatrix:}

\begin{itemize}[leftmargin=*]
    \item \textbf{Nutze Prompt Engineering}, wenn die Aufgabe klar definiert ist und das Modell das Wissen bereits besitzt.
    \item \textbf{Nutze RAG}, wenn faktenbasierte Antworten mit aktuellem Wissen benötigt werden.
    \item \textbf{Nutze Fine-Tuning}, wenn das Modell ein spezifisches Format, einen Stil oder eine Denkweise erlernen muss, die durch Prompts nicht zuverlässig erreichbar ist.
    \item \textbf{Kombiniere alle drei}, wenn das Modell domänenspezifisch antworten \textit{und} auf aktuelle Daten zugreifen soll (Fine-Tuned Base + RAG-Layer).
\end{itemize}

% ===========================================================%
% ARCHITEKTUR
% ===========================================================%
\section{Architektur -- PEFT (Parameter-Efficient Fine-Tuning)}

Full Fine-Tuning -- das Aktualisieren aller Modellgewichte -- ist für LLMs mit 70B+ Parametern praktisch unmöglich ausserhalb von Grosskonzernen. Eine vollständige Fine-Tuning-Runde für Llama 70B kostet auf 8xH100 rund 50.000~USD.

\textbf{PEFT (Parameter-Efficient Fine-Tuning)} aktualisiert nur einen winzigen Bruchteil der Parameter. Der dominierende Ansatz ist \textbf{LoRA (Low-Rank Adaptation)}.

\subsection{Wie LoRA funktioniert}

LoRA fügt für jede trainierbare Schicht zwei kleine Matrizen $A$ und $B$ hinzu, deren Produkt $A \cdot B$ die Gewichtsänderung $\Delta W$ approximiert:

\begin{verbatim}
   Original Forward:   y = W * x
   LoRA Forward:       y = W * x + (B * A) * x
                                \_______/
                            rang-reduziertes Update

   W:    Originalgewichte (eingefroren, nicht trainiert)
   A, B: LoRA-Matrizen (trainiert)
   Rang: Kontrolliert die Anzahl trainierbarer Parameter
\end{verbatim}

Bei einem typischen LLM mit 4.096 Hidden-Dimension und LoRA-Rank~16:

\begin{itemize}[leftmargin=*]
    \item Eine lineare Schicht (4.096 $\times$ 4.096): 16,8 Mio Parameter
    \item LoRA-Matrizen A (4.096 $\times$ 16) + B (16 $\times$ 4.096): 131.000 Parameter
    \item \textbf{Reduktion: Faktor 128}
\end{itemize}

Ein 70B-Modell benötigt mit Full Fine-Tuning 140~GB VRAM. Mit QLoRA (Quantized LoRA) sind 48~GB ausreichend -- eine einzelne A100~80GB genügt.

\subsection{LoRA-Konfiguration}

\begin{lstlisting}[language=Python, caption=LoRA-Konfiguration mit Hugging Face PEFT]
from peft import LoraConfig, get_peft_model

lora_config = LoraConfig(
    r=16,                     # Rang der LoRA-Matrizen
    lora_alpha=32,            # Skalierungsfaktor (alpha / r)
    target_modules=[          # Welche Schichten erhalten LoRA?
        "q_proj",             # Query-Projektion
        "k_proj",             # Key-Projektion
        "v_proj",             # Value-Projektion
        "o_proj",             # Output-Projektion
    ],
    lora_dropout=0.05,        # Dropout für Regularisierung
    bias="none",              # Bias-Gewichte nicht trainieren
    task_type="CAUSAL_LM",    # Kausales Sprachmodell
)

model = AutoModelForCausalLM.from_pretrained(
    "meta-llama/Llama-3.2-8B-Instruct",
    torch_dtype=torch.bfloat16,
    device_map="auto",
)

# LoRA-Adapter an das Modell heften
model = get_peft_model(model, lora_config)
print(f"Trainierbare Parameter: "
      f"{model.num_parameters(only_trainable=True):,} "
      f"von {model.num_parameters():,} "
      f"({model.num_parameters(only_trainable=True) / model.num_parameters():.2%})")
\end{lstlisting}

\subsubsection{Rank-Wahl}

Der Rang $r$ bestimmt die Anzahl trainierbarer Parameter und damit die Kapazität des Adapters:

\begin{table}[H]
\centering
\begin{tabularx}{\linewidth}{lXXX}
\toprule
\textbf{Rank} & \textbf{Trainierbare Param. (8B)} & \textbf{VRAM-Zusatz} & \textbf{Einsatz} \\
\midrule
4 & 2,1 Mio & 8 MB & Einfache Formatanpassung \\
8 & 4,2 Mio & 16 MB & Leichte Stil-Anpassung \\
16 & 8,4 Mio & 32 MB & Standard -- neue Aufgaben \\
32 & 16,8 Mio & 64 MB & Komplexe Domänen (Code, Medizin) \\
64 & 33,6 Mio & 128 MB & Starke Anpassung (mehrere Tasks) \\
\bottomrule
\end{tabularx}
\caption{Rang-Konfiguration und Auswirkungen bei 8B-Modell}
\label{tab:lora-rank}
\end{table}

\textbf{Richtwert:} Für die meisten Aufgaben reicht $r = 16$. Nur wenn das Modell eine komplett neue Fähigkeit erlernen soll (z.\,B. Tool-Use-Format), sind höhere Ränge nötig.

\subsubsection{Target-Module}

Welche Schichten erhalten LoRA? Die Standardwahl ist \texttt{q\_proj, v\_proj}:

\begin{itemize}[leftmargin=*]
    \item \texttt{q\_proj + v\_proj}: Bewährt für die meisten Aufgaben
    \item \texttt{q\_proj + k\_proj + v\_proj + o\_proj}: Mehr Kapazität, kaum zusätzlicher VRAM
    \item \texttt{alle linear layers}: Maximale Kapazität (immer mit niedrigem Rang kombinieren)
\end{itemize}

\subsection{QLoRA -- Quantisiertes Fine-Tuning}

QLoRA kombiniert zwei Techniken: Das Basismodell wird in 4-Bit (NF4) quantisiert geladen, während die LoRA-Adapter in voller Präzision (FP16/BF16) trainiert werden:

\begin{lstlisting}[language=Python, caption=QLoRA -- Fine-Tuning eines quantisierten Modells]
from transformers import BitsAndBytesConfig

bnb_config = BitsAndBytesConfig(
    load_in_4bit=True,
    bnb_4bit_use_double_quant=True,
    bnb_4bit_quant_type="nf4",
    bnb_4bit_compute_dtype=torch.bfloat16,
)

model = AutoModelForCausalLM.from_pretrained(
    "meta-llama/Llama-3.3-70B-Instruct",
    quantization_config=bnb_config,
    device_map="auto",
)

model = get_peft_model(model, lora_config)
# Modell in 4-Bit geladen (~35 GB VRAM fuer 70B)
# LoRA-Adapter in BF16 trainiert (~8 MB zusaetzlich)
\end{lstlisting}

QLoRA reduziert den VRAM-Bedarf drastisch: Ein 70B-Modell benötigt statt 140~GB nur noch \(\sim 38\)~GB und trainiert auf einer einzelnen A100~80GB.

% ===========================================================%
% PRAXIS
% ===========================================================%
\section{Praxis -- Dataset-Erstellung}

Ein Fine-Tuning-Dataset ist das Fundament der gesamten Optimierung. Ein schlechtes Dataset produziert ein schlechtes Modell -- selbst mit perfekter LoRA-Konfiguration.

\subsection{Instruction-Format}

Das Standardformat für Chat-Modelle ist das ShareGPT- oder ChatML-Format:

\begin{lstlisting}[language=Python, caption=Instruction-Dataset im Chat-Format]
DATASET = [
    {
        "messages": [
            {
                "role": "system",
                "content": "Du bist ein SQL-Experte. "
                           "Generiere nur SQL-Queries, keine Erklaerungen.",
            },
            {
                "role": "user",
                "content": "Finde alle Kunden, die in den letzten 30 Tagen "
                           "mehr als 500 EUR ausgegeben haben.",
            },
            {
                "role": "assistant",
                "content": (
                    "SELECT c.customer_id, c.name, SUM(o.total) as total_spent\n"
                    "FROM customers c\n"
                    "JOIN orders o ON c.customer_id = o.customer_id\n"
                    "WHERE o.order_date >= DATE('now', '-30 days')\n"
                    "GROUP BY c.customer_id\n"
                    "HAVING total_spent > 500\n"
                    "ORDER BY total_spent DESC;"
                ),
            },
        ]
    },
    # 500+ weitere Beispiele
]
\end{lstlisting}

\subsection{Dataset-Qualitätskriterien}

\begin{table}[H]
\centering
\begin{tabularx}{\linewidth}{lX}
\toprule
\textbf{Kriterium} & \textbf{Beschreibung} \\
\midrule
Minimale Grösse & 100~Beispiele pro Task, 500+ für zuverlässige Ergebnisse \\
Diversität & Nicht 100x dasselbe Muster -- echte Varianz in Formulierungen \\
Korrektheit & Jede Antwort muss faktengeprüft sein. Ein Fehler trainiert das Modell falsch \\
Helping vs Eval & 80\,\% Instruction-following, 20\,\% Chain-of-Thought für Reasoning \\
Balancing & Keine Kategorie dominiert. Bei Ungleichgewicht: undersamplen \\
Deduplizierung & Exakt gleiche Paare entfernen (Levenshtein-Ähnlichkeit $>$ 0.9) \\
\bottomrule
\end{tabularx}
\caption{Dataset-Qualitätskriterien für Fine-Tuning}
\label{tab:dataset-quality}
\end{table}

\subsection{Dataset-Splits}

\begin{lstlisting}[language=Python, caption=Dataset-Validierung und Splits]
from datasets import Dataset, DatasetDict

# Dataset splitten
dataset = Dataset.from_list(DATASET)
splits = dataset.train_test_split(test_size=0.1, seed=42)

dataset_dict = DatasetDict({
    "train": splits["train"],
    "test": splits["test"],
})

# Qualitaetscheck: Token-Laenge pro Beispiel
from transformers import AutoTokenizer

tokenizer = AutoTokenizer.from_pretrained(
    "meta-llama/Llama-3.2-8B-Instruct"
)

def token_length(example):
    # Nur Assistant-Tokens zaehlen (Output-Learning)
    assistant_text = [
        msg["content"] for msg in example["messages"]
        if msg["role"] == "assistant"
    ][0]
    tokens = tokenizer(assistant_text, truncation=True)
    example["assistant_tokens"] = len(tokens["input_ids"])
    return example

dataset = dataset.map(token_length)

# Verteilung der Output-Laengen prüfen
lengths = dataset["assistant_tokens"]
print(f"Min: {min(lengths)}, Max: {max(lengths)}, "
      f"Avg: {sum(lengths)/len(lengths):.0f}")

# Warnung bei zu kurzen Outputs
if sum(1 for l in lengths if l < 10) > len(lengths) * 0.1:
    print("WARN: Ueber 10% der Outputs haben <10 Tokens")
\end{lstlisting}

% ===========================================================%
% PERFORMANCE
% ===========================================================%
\section{Training und Evaluation}

\subsection{Trainings-Infrastruktur}

\begin{table}[H]
\centering
\begin{tabularx}{\linewidth}{lXXX}
\toprule
\textbf{Setup} & \textbf{GPU} & \textbf{Batch/Step} & \textbf{Kosten pro Stunde} \\
\midrule
8B LoRA & 1x A100~80GB & 32 & 2--3 \$ \\
8B QLoRA & 1x A10~24GB & 8 & 0,50--1 \$ \\
70B QLoRA & 1x A100~80GB & 8 & 2--3 \$ \\
70B Full FT & 8x H100 & 32 & 40--60 \$ \\
\bottomrule
\end{tabularx}
\caption{Trainings-Infrastruktur: GPU-Anforderungen und Kosten}
\label{tab:train-setup}
\end{table}

\begin{lstlisting}[language=Python, caption=Fine-Tuning mit Hugging Face Trainer]
from transformers import TrainingArguments, Trainer

training_args = TrainingArguments(
    output_dir="./lora-llama-sql",
    per_device_train_batch_size=8,
    gradient_accumulation_steps=4,   # Effektive Batch-Size: 32
    learning_rate=2e-4,
    warmup_steps=100,
    num_train_epochs=3,
    logging_steps=10,
    save_strategy="epoch",
    evaluation_strategy="epoch",
    fp16=True,
    report_to="wandb",              # Experiment-Tracking
)

trainer = Trainer(
    model=model,
    args=training_args,
    train_dataset=dataset_dict["train"],
    eval_dataset=dataset_dict["test"],
    tokenizer=tokenizer,
    data_collator=lambda data: tokenizer.pad(
        [{"text": format_chat(d["messages"])} for d in data],
        return_tensors="pt",
    ),
)

trainer.train()
\end{lstlisting}

\subsection{Evaluation nach dem Training}

Nach dem Fine-Tuning muss das Modell gegen zwei Metriken evaluiert werden:

\begin{description}[leftmargin=*,style=nextline]
    \item[Task-Performance] Wie gut ist das Modell in der \textit{neuen} Aufgabe? Gemessen gegen das Test-Set.
    \item[Catastrophic Forgetting] Wie stark hat das Modell \textit{alte} Fähigkeiten verloren? Gemessen gegen ein General-Benchmark (MMLU, HumanEval).
\end{description}

\begin{lstlisting}[language=Python, caption=Forgetting-Check nach Fine-Tuning]
BENCHMARK = [
    {"input": "Was ist die Hauptstadt von Frankreich?",
     "expected": "Paris"},
    {"input": "Berechne 25 * 4 + 100",
     "expected": "200"},
    {"input": "Erklaere den Begriff 'RAG' in einem Satz.",
     "expected_contains": ["Retrieval", "Augmented", "Generation"]},
]

def evaluate_forgetting(model, tokenizer, benchmark):
    before_scores = run_benchmark("baseline_model", benchmark)
    after_scores = run_benchmark(model, benchmark)

    for task, before, after in zip(benchmark, before_scores, after_scores):
        diff = after - before
        status = "OK" if diff >= -0.05 else "FORGETTING"
        print(f"{status}: {task['input'][:40]}... "
              f"{before:.0%} -> {after:.0%} ({diff:+.0%})")

    # Gesamtdrift
    avg_drift = sum(
        a - b for a, b in zip(after_scores, before_scores)
    ) / len(benchmark)
    print(f"\nDurchschnittlicher Drift: {avg_drift:+.0%}")

    if avg_drift < -0.10:
        print("WARN: Staerker Drift! Dataset hat Allgemeinwissen beschaedigt.")
    elif avg_drift < -0.05:
        print("HINWEIS: Leichter Drift erkennbar. LoRA-Rang oder "
              "Lernrate pruefen.")
    else:
        print("OK: Kein signifikanter Knowledge-Drift.")
\end{lstlisting}

\subsection{Multi-LoRA-Serving}

In Produktion willst du oft mehrere LoRA-Adapter parallel bereitstellen (ein Adapter für SQL, einer für E-Mail, einer für Code). vLLM unterstützt Multi-LoRA-Serving nativ:

\begin{lstlisting}[language=bash, caption=Multi-LoRA-Serving mit vLLM]
# Server starten mit mehreren LoRA-Adaptern
vllm serve meta-llama/Llama-3.3-70B-Instruct \
    --lora-modules sql=/path/to/lora-sql \
    --lora-modules email=/path/to/lora-email \
    --lora-modules code=/path/to/lora-code \
    --max-lora-rank 32 \
    --enable-lora

# Client: Adapter per Header waehlen
response = client.chat.completions.create(
    model="sql",  # LoRA-Adapter-Name
    messages=[{"role": "user", "content": "SELECT * FROM users..."}],
    extra_body={"max_tokens": 256},
)
\end{lstlisting}

Der Adapter wird bei jedem Request dynamisch in den Forward-Pass eingeblendet -- ohne Modell-Neuladen. Der Speicher-Overhead pro Adapter beträgt nur einige MB.

% ===========================================================%
% BEST PRACTICES
% ===========================================================%
\section{Best Practices}

\begin{itemize}[leftmargin=*]
    \item \textbf{Lernrate}: LoRA benötigt höhere Lernraten (1e-4 bis 5e-4) als Full Fine-Tuning (1e-5). Starte mit 2e-4.
    \item \textbf{Rang nicht überdimensionieren}: $r=16$ deckt 90\,\% der Aufgaben ab. Höherer Rang = mehr Overfitting-Risiko.
    \item \textbf{Target-Module}: q\_proj + v\_proj + o\_proj geben die beste Abdeckung bei geringem Overhead.
    \item \textbf{Dataset-Qualität vor Quantität}: 500 hochwertige Beispiele sind besser als 5.000 automatisierte.
    \item \textbf{Learning Rate Warmup}: 100--200 Schritte Warmup verhindern anfängliche Destabilisierung.
    \item \textbf{Adapter mergen}: Nach dem Training kann der LoRA-Adapter in die Originalgewichte gemergt werden (kein separater Adapter nötig): \texttt{model = model.merge\_and\_unload()}.
    \item \textbf{Checkpointing}: Speichere alle 500 Schritte. Bei Overfitting zum früheren Checkpoint zurückkehren.
\end{itemize}

% ===========================================================%
% ANTI-PATTERNS
% ===========================================================%
\section{Anti-Patterns}

\begin{itemize}[leftmargin=*]
    \item \textbf{Auf synthetische Daten allein vertrauen}: Ein vollständig synthetisches Dataset übernimmt die Fehler des Generierungsmodells. Immer 20\,\% manuell geprüfte Beispiele.
    \item \textbf{Kein Forgetting-Check}: Ein FT-Modell, das auf SQL glänzt, aber keine einfache Frage mehr beantworten kann, ist unbrauchbar. Immer General-Benchmarks mitlaufen lassen.
    \item \textbf{Zu viele Epochen}: Nach 3--5 Epochen beginnt Overfitting. Früher Stop anhand des Test-Verlusts.
    \item \textbf{Lange Outputs im Dataset}: Wenn die Assistant-Antworten durchschnittlich 2.000 Tokens lang sind, trainiert das Modell hauptsächlich, wie man lang antwortet -- nicht wie man richtig antwortet.
    \item \textbf{FT des gesamten Modells ohne Not}: Full Fine-Tuning ist 100--1000x teurer als LoRA und erhöht das Overfitting-Risiko massiv. Nur bei grundlegend neuen Fähigkeiten notwendig.
\end{itemize}

% ===========================================================%
% ZUSAMMENFASSUNG
% ===========================================================%
\begin{zusammenfassung}
\begin{itemize}[leftmargin=*]
    \item Fine-Tuning optimiert Modellgewichte für spezifische Aufgaben, während RAG den Prompt mit Kontext anreichert -- beide ergänzen sich
    \item LoRA trainiert rang-reduzierte Adapter-Matrizen statt aller Gewichte: 99,9\,\% Parameterreduktion
    \item QLoRA ermöglicht Fine-Tuning von 70B-Modellen auf einer einzelnen GPU (quantisiertes Basismodell + LoRA in FP16)
    \item Das Dataset ist der wichtigste Erfolgsfaktor -- 500 qualitätsgeprüfte Beispiele schlagen 5.000 synthetische
    \item Forgetting-Check nach jedem FT ist Pflicht: Das Modell darf alte Fähigkeiten nicht verlieren
    \item Multi-LoRA-Serving (vLLM) erlaubt dozens von Adaptern parallel auf einem Modell
\end{itemize}
\end{zusammenfassung}

\begin{merke}
Fine-Tuning ist kein Ersatz für RAG, sondern eine Ergänzung. RAG liefert Fakten, FT liefert Format und Stil. Die beste Architektur kombiniert ein fine-getuntes Basismodell mit einem RAG-Layer für aktuelle Informationen.
\end{merke}

% ===========================================================%
% PRAXISPROJEKT
% ===========================================================%
\section{Praxisprojekt -- SQL-Experte mit LoRA}

\textbf{Ziel:} Fine-Tune ein 8B-LLM auf SQL-Generierung mit LoRA.

\textbf{Aufgaben:}
\begin{enumerate}[leftmargin=*]
    \item Erstelle ein Dataset mit 300 SQL-Query-Paaren (vom Typ "Frage $\rightarrow$ SQL")
    \item Konfiguriere LoRA mit $r=16$, target\_modules = q\_proj, v\_proj, o\_proj
    \item Führe das Training für 3 Epochen durch
    \item Validiere: Test-Accuracy und Forgetting-Check gegen MMLU-Stichprobe
    \item Deploye den Adapter mit vLLM Multi-LoRA-Serving
\end{enumerate}

\textbf{Erfolgskriterium:} Das Modell generiert korrektes SQL für 90\,\% der Test-Query. Die MMLU-Accuracy sinkt um maximal 5\,\%.

% ===========================================================%
% INTERVIEWFRAGEN
% ===========================================================%
\section{Interviewfragen}

\begin{enumerate}[leftmargin=*]
    \item \textbf{Wann würdest du Fine-Tuning statt RAG einsetzen?}

    \textbf{Antwort:} Fine-Tuning, wenn das Modell ein spezifisches Output-Format erlernen muss (SQL, JSON-Schema, E-Mail-Stil), das durch Prompting nicht zuverlässig erreichbar ist. RAG, wenn faktenbasierte Antworten mit aktuellem Wissen benötigt werden. Kombination: FT für Format/Kompetenz + RAG für Fakten.

    \item \textbf{Erkläre den Mechanismus von LoRA. Was ist der Rang?}

    \textbf{Antwort:} LoRA fügt zwei trainierbare Matrizen A und B pro Schicht hinzu, deren Produkt die Gewichtsänderung approximiert. Rang r bestimmt die Dimension der Matrizen: A hat Shape (hidden, r), B hat Shape (r, hidden). Höherer Rang = mehr Kapazität, aber auch mehr Overfitting.

    \item \textbf{Dein fine-getuntes Modell ist super in der neuen Aufgabe, beantwortet aber keine Allgemeinwissen-Fragen mehr. Was ist passiert?}

    \textbf{Antwort:} Catastrophic Forgetting. Das Dataset enthielt zu wenige allgemeine Beispiele oder die Lernrate war zu hoch. Lösung: Dataset um 20\,\% allgemeine Wissensfragen erweitern, Lernrate senken, LoRA-Rang prüfen (niedriger = weniger Änderung).

    \item \textbf{Wie gehst du vor, wenn du LoRA-Gewichte in Produktion bringen willst, aber nur einen Inference-Server hast?}

    \textbf{Antwort:} Multi-LoRA-Serving mit vLLM: Der Server lädt das Basismodell einmal und blendet pro Request den passenden Adapter ein. Speicher-Overhead pro Adapter: einige MB. alternative: LoRA-Gewichte in das Basismodell mergen und deployen (kein separater Adapter nötig).

    \item \textbf{Wie gross sollte ein Fine-Tuning-Dataset mindestens sein?}

    \textbf{Antwort:} 100 Beispiele pro Task, 500+ für zuverlässige Ergebnisse. Wichtiger als Quantität ist Qualität: Ein fehlerhaftes Beispiel trainiert das Modell falsch. 500 manuell geprüfte Beispiele sind wertvoller als 5.000 automatisierte.

    \item \textbf{Welche Lernrate wählst du für LoRA und warum?}

    \textbf{Antwort:} 2e-4 als Startwert. LoRA benötigt höhere Lernraten als Full Fine-Tuning (1e-5), weil nur wenige Parameter trainiert werden. Bei Lernraten unter 1e-4 passiert kaum etwas, über 5e-4 droht Instabilität.

    \item \textbf{Du hast nur eine T4 GPU (16 GB). Welches Fine-Tuning ist möglich?}

    \textbf{Antwort:} QLoRA mit einem 8B-Modell in 4-Bit: 16~GB reichen für Training (8B ~ 5~GB in 4-Bit + Adapter + Optimizer-States). Batch-Size auf 1--2 reduzieren, Gradient Accumulation nutzen. 70B ist auf einer T4 nicht möglich.
\end{enumerate}

% ===========================================================%
% RESSOURCEN
% ===========================================================%
\section{Ressourcen}

\begin{itemize}[leftmargin=*]
    \item \textbf{Hugging Face PEFT}: \href{https://github.com/huggingface/peft}{github.com/huggingface/peft} -- LoRA/QLoRA/DoRA Bibliothek
    \item \textbf{Axolotl}: \href{https://github.com/axolotl-ai-cloud/axolotl}{github.com/axolotl-ai-cloud/axolotl} -- Feinabgestimmtes FT-Framework
    \item \textbf{Unsloth}: \href{https://github.com/unslothai/unsloth}{github.com/unslothai/unsloth} -- 2x schnelleres FT mit optimiertem Kernel
    \item \textbf{OpenAI Fine-Tuning API}: \href{https://platform.openai.com/docs/guides/fine-tuning}{platform.openai.com} -- Managed Fine-Tuning
    \item \textbf{QLoRA Paper}: Dettmers et al. (2023): \textit{QLoRA: Efficient Finetuning of Quantized Language Models}
\end{itemize}

\textbf{Nächster Schritt:} Im nächsten Kapitel erfährst du, wie du LLM-Systeme durch Caching, intelligentes Routing und Guardrails optimierst und absicherst.
\endinput
